\section{Pilot Results}
\subsection{Integrity outcome}
P001 completed all ten recorded integrity checks. Table~\ref{tab:integrity} is generated
from the immutable raw record; the overall gate passed while the record remains marked
pilot and not publication-eligible.

This outcome is procedural rather than substantive. It establishes that the recorded run
contains the required manifests, integrity checks, metadata, and metric outputs in a form
that can be traced through the analysis pipeline. It does not validate the sensitivity of a
particular deletion assay, the competitive quality of an approximate method, or a
counterfactual-equivalence claim. Keeping these interpretations separate is essential: an
auditable run can still be scientifically uninformative for the question it was intended to
test.
\begin{table}[ht]
\centering
\caption{P001 integrity-gate outcomes generated from the raw experiment record.}
\label{tab:integrity}
\small
\begin{tabular}{ll}
\toprule
Gate & P001 status \\
\midrule
All Numeric Outputs Finite & PASS \\
Checkpoints Present And Hashed & PASS \\
Configuration Hash Present & PASS \\
Exact Exclusion Assertion & PASS \\
Forget Retain Disjoint & PASS \\
Multiple Exact Seeds & PASS \\
Partition Reconstructs Train & PASS \\
Process Completed & PASS \\
Raw Predictions And Labels Present & PASS \\
Recognized Maturity & PASS \\
Runtime Complete & PASS \\
\bottomrule
\end{tabular}

\end{table}

\subsection{Exact-retraining population}
The 15 exact replicas exhibit nonzero variability on all reported splits. Table
~\ref{tab:exact-ranges} reports per-seed ranges and means. Forget-set accuracy ranges from
0.8222 to 0.9222, while test accuracy ranges from 0.8096 to 0.8504. These ranges illustrate
why neither one exact model nor one absolute forget-set value is an adequate target.

The displayed ranges are descriptive summaries of this finite ensemble, not confidence
intervals for a population of all possible trainings. Their main role is to make the
comparison target explicit. An approximate model should be considered in relation to the
dispersion among models trained on the retained dataset, rather than treated as successful
or unsuccessful simply because it differs from one selected retraining seed. Later studies
will need to report uncertainty for the envelope itself and assess whether its estimates
are stable under additional seeds and independent original models.
\begin{table}[ht]
\centering
\caption{Metric ranges across 15 independently exact-retrained models. NLL is negative
log-likelihood and ECE uses 15 bins.}
\label{tab:exact-ranges}
\small
\begin{tabular}{lrrr}
\toprule
Split and metric & Minimum & Mean & Maximum \\
\midrule
Forget Accuracy & 0.8222 & 0.8741 & 0.9222 \\
Forget NLL & 0.3426 & 0.4040 & 0.4750 \\
Forget Brier & 0.1568 & 0.1959 & 0.2356 \\
Forget ECE & 0.0510 & 0.0839 & 0.1240 \\
Retain Accuracy & 0.8263 & 0.8509 & 0.8695 \\
Retain NLL & 0.3579 & 0.4078 & 0.4587 \\
Retain Brier & 0.1861 & 0.2128 & 0.2410 \\
Retain ECE & 0.0086 & 0.0248 & 0.0447 \\
Test Accuracy & 0.8096 & 0.8373 & 0.8504 \\
Test NLL & 0.4145 & 0.4584 & 0.5111 \\
Test Brier & 0.2116 & 0.2333 & 0.2628 \\
Test ECE & 0.0120 & 0.0204 & 0.0422 \\
\bottomrule
\end{tabular}

\end{table}

\subsection{Frozen seed-count diagnostic}
The frozen $k=10\rightarrow15$ diagnostics changed by \PilotJSConvergence{} for mean
Jensen--Shannon divergence and \PilotDisagreementConvergence{} for prediction disagreement.
Both values fall below the predefined 10\% pilot threshold. The ensemble therefore passed
the stated pilot convergence diagnostic for these two metrics. This does not establish
that 15 seeds are sufficient for other datasets, architectures, deletion regimes, or
metrics, and the threshold is not an equivalence margin.

The diagnostic was intentionally limited to two inexpensive predictive distances. It asks
whether enlarging the finite reference ensemble changed those descriptive means by less
than a predeclared pilot tolerance; it does not demonstrate convergence of all fidelity
dimensions. In particular, calibration, local behavior, membership-inference outcomes,
and representation measures may require different seed counts. The procedure also did not
adaptively stop once favorable values were observed, which preserves its role as a fixed
engineering check rather than an outcome-driven sample-size decision.

\subsection{Approximate-method descriptives}
Table~\ref{tab:methods} records the executed method runtimes and accuracies. The values
confirm that the common method and metric interfaces executed, but they are not a ranking.
In particular, high accuracy on the IID forget set does not demonstrate that the deleted
records remain memorized: exact models that never trained on those records also achieve
substantial forget-set accuracy.

This distinction is especially consequential for semantically ordinary IID examples. A
counterfactual model can classify a deleted garment correctly because it learned related
structure from retained examples, while an unlearned model can reduce accuracy through
collateral damage rather than faithful removal of the deleted example's influence. The
descriptive table should therefore be read as evidence that the baseline interfaces ran and
produced comparable artifacts, not as evidence that one baseline is more effective than
another.
\begin{table}[ht]
\centering
\caption{Descriptive P001 outputs for approximate methods. These values are not interpreted
as confirmatory unlearning outcomes because the deletion assay was retrospectively U0.}
\label{tab:methods}
\small
\begin{tabular}{lrrrr}
\toprule
Method & Seconds & Forget acc. & Retain acc. & Test acc. \\
\midrule
Fine-tuning & 0.27 & 0.900 & 0.865 & 0.845 \\
Naive gradient ascent & 1.15 & 0.900 & 0.863 & 0.849 \\
SCRUB & 1.87 & 0.922 & 0.869 & 0.850 \\
\bottomrule
\end{tabular}

\end{table}

\subsection{Retrospective identifiability analysis}
A retrospective separability analysis compared the original model with the exact
ensemble. The separability ratios were \PilotJSSeparability{} for JS divergence and
\PilotDisagreementSeparability{} for prediction disagreement. Both were below their
exploratory power-derived thresholds, yielding U0 (assay-insensitive). Consequently, the
pilot does not determine whether any approximate method forgot the IID-deleted examples.
No confirmatory method result is reported.

U0 is an assay statement, not a statement that deleted examples had no effect on training.
It means that, under the archived perturbation construction and selected distances, the
original-versus-exact contrast was not sufficiently separated to support the intended
decision. A stronger deletion regime, a different probe distribution, a longer training
schedule, or a different architecture may yield a measurable contrast; those are
hypotheses for prospectively specified studies rather than explanations that can rescue the
current pilot.

\subsection{Prospective class-conditioned assay calibration}
A separate prospective pilot used five original seeds and 15 exact-retraining seeds
per condition, without evaluating approximate methods. Full deletion of Fashion-MNIST
class 0 and class 5 produced JS-divergence separability ratios of
\PhaseThreeClassZeroJSRatio{} and \PhaseThreeClassFiveJSRatio{}, exceeding their
simulated minimum detectable ratios of \PhaseThreeClassZeroMinimum{} and
\PhaseThreeClassFiveMinimum{}. Half-class deletion and every tested disagreement,
CKA-distance, and MMD endpoint remained assay-insensitive. A model-level precision
simulation selected ten exact seeds for JS under the frozen 20\% interval rule in both
full-class conditions.

\subsection{Prospective equivalence-margin calibration}
A second frozen pilot tested whether the condition-specific exact/exact 95th percentile
could detect a controlled violation in which 1\% of the forbidden class remained in
training. The candidate margins were \PhaseThreeClassZeroMargin{} for class 0 and
\PhaseThreeClassFiveMargin{} for class 5. Detection probabilities were only
\PhaseThreeClassZeroDetection{} and \PhaseThreeClassFiveDetection{}, respectively,
below the prospectively required 80\% in both cases. The candidate margins are therefore
inadmissible. Larger contamination levels were designated diagnostics and cannot rescue
the failed primary rule. No approximate-method equivalence claim is made.


\subsection{Prospective v2.1 calibration}
The replacement class-stratified calibration was frozen before the holdout exact models were
trained. It uses a validation-only positive control that swaps the top two probabilities on
one example in each class of the fixed 1,000-example probe. Fifteen calibration seeds estimate
each condition's one-sided lower-bound excess-JS margin; ten disjoint heldout seeds test the
predeclared 80\% exceedance and monotonic-severity requirements. Both class-deletion conditions
passed: the candidate margins were \VTwoOneClassZeroMargin{} for class 0 and
\VTwoOneClassFiveMargin{} for class 5, with a heldout primary exceedance rate of
\VTwoOneHoldoutRate{} and strictly increasing median effects at 0.5\%, 1\%, and 2\% severity.

This is validation-only evidence that the selected JS assay can detect and calibrate an
interpretable, class-balanced synthetic perturbation in this Fashion-MNIST/small-CNN setting.
It does not show that any approximate method falls within either margin, establish a general
equivalence threshold, rescue the failed v1 calibration, or support a certified-removal claim.
The margin may be used only by a separately frozen method comparison that preserves this
dataset, architecture, probe, and endpoint scope.

\subsection{Runtime}
Table~\ref{tab:runtime} separates original training, aggregate exact retraining, and the
aggregate approximate-method execution. Runtime values characterize this software and
hardware configuration and are not universal performance measurements. Speed alone has no
scientific meaning without an interpretable fidelity comparison.

The recorded timing is deliberately narrow: it measures the wall-clock scope defined by the
runner on one ROCm environment. It is not an estimate of energy consumption, peak memory,
cloud cost, or accelerator-independent throughput. A later compute--fidelity frontier must
standardize timing boundaries, preserve device and driver metadata, and report both the
cost of the approximate update and the cost of constructing the exact-retraining reference
population needed to evaluate it.
\begin{table}[ht]
\centering
\caption{Recorded P001 runtime on the RX 7800 XT/ROCm environment.}
\label{tab:runtime}
\begin{tabular}{lr}
\toprule
Operation & Seconds \\
\midrule
Original training & 9.59 \\
15 exact retrainings, aggregate & 17.46 \\
Approximate methods, aggregate & 3.29 \\
\bottomrule
\end{tabular}

\end{table}
